% Generated draft: requires mathematical and source review.
\documentclass[11pt]{article}
\usepackage[margin=25mm]{geometry}
\usepackage{fontspec}
\setmainfont{Latin Modern Roman}
\usepackage{amsmath,amssymb,mathtools,unicode-math}
\setmathfont{Latin Modern Math}
\usepackage{enumitem,needspace,xurl,hyperref}
\setlength{\parindent}{0pt}
\setlength{\parskip}{5pt}
\setlength{\emergencystretch}{3em}
\newcommand{\N}{\mathbb{N}}
\newcommand{\Z}{\mathbb{Z}}
\newcommand{\Q}{\mathbb{Q}}
\newcommand{\R}{\mathbb{R}}
\newcommand{\C}{\mathbb{C}}
\newcommand{\F}{\mathbb{F}}
\newcommand{\A}{\mathbb{A}}
\newcommand{\PP}{\mathbb{P}}
\newcommand{\E}{\mathbb{E}}
\newcommand{\eps}{\varepsilon}
\newcommand{\dd}{\,\mathrm{d}}
\begin{document}
\section*{UnsolvedMath ID: 701347}
\textbf{Source problem code:} catalogue entry
\subsection{Definitions and mathematical statement}
Let \(G\) be a group. For \(H\leq G\) and \(g\in G\), write \(H^g=g^{-1}Hg\), and write \(\langle A,B\rangle\) for the subgroup generated by subgroups \(A\) and \(B\). A subgroup \(H\leq G\) is \emph{pronormal} if
\[
\forall g\in G\quad \exists x\in\langle H,H^g\rangle\quad H^x=H^g.
\]
It is \emph{strongly pronormal} if
\[
\forall L\leq H\quad\forall g\in G\quad
\exists x\in\langle H,L^g\rangle\quad (L^g)^x\leq H.
\]
For finite \(G\), a subgroup \(H\leq G\) is a \emph{Hall subgroup} if
\[
\gcd\bigl(|H|,[G:H]\bigr)=1.
\]
Here, having a \emph{Sylow tower} is interpreted as follows: if \(p_1,\ldots,p_k\) are the distinct prime divisors of \(|H|\), there is some ordering of these primes and a series
\[
1=H_0\leq H_1\leq\cdots\leq H_k=H,
\qquad H_i\trianglelefteq H\quad(0\leq i\leq k),
\]
such that \(|H_i/H_{i-1}|\) is the largest power of \(p_i\) dividing \(|H|\), for every \(1\leq i\leq k\). The trivial group has the empty tower, with \(k=0\). The supplied source does not expand the term ``Sylow tower'' or specify an ordering convention; the interpretation above requires confirmation against the source.

\textbf{Remaining question, part (c).} Is it true that for every finite group \(G\) and every Hall subgroup \(H\leq G\) having a Sylow tower, \(H\) is strongly pronormal in \(G\)? Explicitly, under these hypotheses, must
\[
\forall L\leq H\quad\forall g\in G\quad
\exists x\in\langle H,L^g\rangle\quad (L^g)^x\leq H
\]
hold? No simplicity assumption is imposed on \(G\).

\textbf{Archived questions and supplied qualifications.} Part (a) asks whether every Hall subgroup of every finite simple group is pronormal. Part (b) asks whether every Hall subgroup of every finite simple group is strongly pronormal. The supplied Archive of Solved Problems entry reports an affirmative answer to (a), attributed to E. P. Vdovin and D. O. Revin (2012), and a negative answer to (b), attributed to M. N. Nesterov (2017). The source additionally reports that an affirmative answer to (a) for Hall subgroups of odd order was known modulo the classification of finite simple groups (CFSG), citing F. Gross (1987). It also states that some finite nonsimple groups have nonpronormal Hall subgroups, whereas Hall subgroups having a Sylow tower are known to be pronormal. These are preserved source reports, not independently verified mathematical or literature assessments. The registry labels the whole record partially solved and places part (c) among the remaining questions.
\subsection{Short English statement}
Part (c) asks whether every Hall subgroup with a Sylow tower in a finite group is strongly pronormal: for every subgroup \(L\leq H\) and every \(g\in G\), can \(L^g\) be conjugated into \(H\) using an element of \(\langle H,L^g\rangle\)? The supplied record reports parts (a) and (b) as resolved. The precise Sylow tower convention needs source confirmation.
\subsection{Sources}
\begin{itemize}
\item Local UnsolvedMath registry: permanent ID \(701347\), exact source problem code \texttt{catalogue entry}, title ``Kourovka Notebook Problem 17.45''; supplied upstream revision \texttt{d371ee917ac9647fbb06cd6fc529deec6f3dda82}. The supplied status is \texttt{partially\_solved}.
\item The Kourovka Notebook (Unsolved Problems in Group Theory), problem 17.45, proposed by E. P. Vdovin and D. O. Revin; 17th issue (2010), with an update attributed in the supplied record to the 21st issue (2026), arXiv:1401.0300v48, \url{https://arxiv.org/abs/1401.0300v48}. The supplied collection description names E. I. Khukhro and V. D. Mazurov as editors. This reference and its version information have not been independently checked.
\item Supplied Archive of Solved Problems reference for part (a): E. P. Vdovin and D. O. Revin, \emph{Siberian Math. J.}, \textbf{53}, no. 3 (2012), 419--430. The archive reports an affirmative answer. No article title or article link is supplied.
\item Supplied Archive of Solved Problems reference for part (b): M. N. Nesterov, \emph{Siberian Math. J.}, \textbf{58}, no. 1 (2017), 128--133. The archive reports that the answer is not always affirmative. No article title or article link is supplied.
\item Supplied background reference concerning odd-order Hall subgroups in part (a): F. Gross, \emph{Bull. London Math. Soc.}, \textbf{19}, no. 4 (1987), 311--319. The supplied qualification is that the affirmative result is known modulo CFSG. No article title or article link is supplied.
\item The supplied status review is dated \texttt{2026-10-16}; its evidence field is ``published'', and it reports parts (a) and (b) as archived solved parts, with the remaining part in the main section. This is supplied dated curation, not an independently verified assessment.
\end{itemize}
\end{document}
